% Ad Atticum 7.12 (ad-atticum-07-12) --- English translation
% Translated by Claude (Anthropic) from the Latin (Perseus).
% Style guide: STYLE.md.

\ciceroLetterOpener{CICERO ATTICO salutem}

\ciceroSection{1} So far I have had one letter from you, sent on the twelfth day before the Kalends, in which there was a hint that you had sent another before it which I had not received. But I beg you to write as often as you can --- not only if you know or have heard anything, but even if you suspect anything; and above all what you think we ought to do, or not to do.

\ciceroSection{2} As for your asking me to take care to let you know what Pompey is doing --- I do not think he himself knows; certainly none of us does. I saw Lentulus the consul at Formiae on the tenth day before the Kalends; I saw Libo: everything is full of fear and going wrong. He is on the way to Larinum; for it is there that the cohorts are, and at Luceria and Teanum, and the rest of them in Apulia. Whether from there he wants to make a stand anywhere or to cross the sea is unknown. If he stays, I fear he cannot have a firm army; if he withdraws, where to or by what route, or what we ought to do, I do not know. For as for that other man, whose \textit{[Greek: Phalarid-tyranny]} you fear, I think he will do everything in the foulest way. Neither the postponement of business nor the withdrawal of the senate and magistrates nor the locking of the treasury will hold him back.

\ciceroSection{3} But these things, as you write, we shall soon know. In the meantime forgive me, please, for writing so much to you so often. For I find rest in it, and I want to draw out letters from you, and above all counsel as to what I should do or in what fashion to bear myself. Shall I throw myself wholly into the cause? I am not deterred by the danger, but I am torn apart by grief. That everything has been done with no plan, or so contrary to my own plan! Or shall I hesitate and turn aside, and give myself up to those who hold the power, who are in possession? \textit{[Greek: I shrink from the Trojans]} --- and not only by the duty of a citizen but by that of a friend I am held back; even so I am often broken by pity for the boys.

\ciceroSection{4} As I, then, am in such disorder --- although the same anxieties grip you --- write something, and above all, if Pompey leaves Italy, what you think we ought to do. Manius Lepidus, for one (we were together), sets that as the limit; Lucius Torquatus the same. As for me, many things, but the lictors not least, get in my way. I have never seen anything that could less be unravelled. So I am asking nothing certain of you yet --- only what you make of it. In a word, your own \textit{[Greek: perplexity]} is what I want to know. That Labienus has left Caesar is now nearly settled.

\ciceroSection{5} If only it had so fallen out that, on coming to Rome, he had found magistrates and a senate in Rome, it would have been of great service to our cause. For it would have looked as though he had condemned the man's wickedness for the commonwealth's sake --- which even now is how it looks, but does less good. For he has no one to do good to, and I think he regrets it; unless, indeed, the news itself is false, that he has left him. For our part, we held it for certain.

\ciceroSection{6} And I should like you, although as you write you are keeping within your own domestic bounds, to set out for me the shape of the City --- whether there is any sense of missing Pompey, any visible hatred of Caesar; and also what you advise about Terentia and Tullia, whether they should be at Rome, or with me, or in some safe place. These things, and anything else, write to me --- or rather, keep writing.
